% Professional Back Cover Page
\clearpage
\newgeometry{margin=0pt}
\thispagestyle{empty}
\pagecolor{coverbottom}
\begin{tikzpicture}[remember picture,overlay]
    % Gradient background (matching front cover)
    \shade[top color=coverbottom,bottom color=covertop]
        (current page.north west) rectangle (current page.south east);

    % Decorative top border
    \fill[accentgold,opacity=0.9]
        ([yshift=-2cm]current page.north west) rectangle ([yshift=-2.2cm]current page.north east);
    \fill[accentlight,opacity=0.7]
        ([yshift=-2.2cm]current page.north west) rectangle ([yshift=-2.25cm]current page.north east);

    % Decorative bottom border
    \fill[accentlight,opacity=0.7]
        ([yshift=2.25cm]current page.south west) rectangle ([yshift=2.2cm]current page.south east);
    \fill[accentgold,opacity=0.9]
        ([yshift=2.2cm]current page.south west) rectangle ([yshift=2cm]current page.south east);

    % Central decorative frame
    \draw[accentgold,line width=2pt,opacity=0.8]
        ([shift={(2cm,2cm)}]current page.south west) rectangle ([shift={(-2cm,-2cm)}]current page.north east);
    \draw[white,line width=0.5pt,opacity=0.6]
        ([shift={(2.1cm,2.1cm)}]current page.south west) rectangle ([shift={(-2.1cm,-2.1cm)}]current page.north east);

    % Content overlay - positioned at page center
    \node[text width=0.6\paperwidth, align=left, anchor=north] at ([yshift=-3.5cm]current page.north) {
        \begin{center}
        {\footnotesize\bfseries\color{white}About This Handbook\par}
        \end{center}

        \vspace{0.15cm}

        {\scriptsize\color{white}
        This comprehensive handbook bridges the gap between high-level TypeScript programming and low-level LLVM IR code generation. The main goal of this project is to create a TypeScript-supported LLVM generator designed for computationally intensive applications, including heterogeneous programming and GPU-accelerated computing. This project will be integrated into more advanced packages of the @euriklis family, providing a robust way for heterogeneous programming via TypeScript. Whether you're building compilers, optimizing performance-critical applications, or targeting multiple architectures, this guide provides practical knowledge for leveraging LLVM's powerful infrastructure.
        \par}

        \vspace{0.2cm}

        {\scriptsize\bfseries\color{accentgold}Main topics:\par}

        \vspace{0.1cm}

        {\tiny\color{white}
        \begin{itemize}[leftmargin=1cm,itemsep=0.03cm]
            \item Core concepts of LLVM IR and its type system
            \item Cross-platform compilation for ARM, x86, and GPU architectures
            \item GPU programming with NVIDIA NVVM, AMD, and SPIR-V
            \item Memory management and optimization techniques
            \item Complete instruction reference and practical examples
        \end{itemize}
        \par}

        \vspace{0.2cm}

        {\scriptsize\bfseries\color{white}About the Author\par}

        \vspace{0.1cm}

        {\tiny\color{white}
        Velislav Karastoychev is an economist and the founder of Euriklis LTD and co-founder of One Agent LTD. With rich experience in econometrics and the implementation of statistical test hypothesis algorithms, he has translated and implemented in TypeScript a vast collection of algorithms for mathematical programming, optimization, linear algebra, and graph/network analysis. His goal is to elevate TypeScript to the level of a comprehensive language for data processing, analysis, and machine learning, bringing the power of low-level optimization to high-level application development.
        \par}

        \vspace{0.2cm}

        \begin{center}
        {\scriptsize\color{white}
        \textbf{GitHub Repository:}\\[0.1cm]
        \url{https://github.com/VelislavKarastoychev/euriklis-llvm-ir}
        \par}

        \vspace{0.15cm}

        {\tiny\color{white!80}
        \textcopyright\ \the\year\ Velislav Karastoychev. All rights reserved.
        \par}
        \end{center}
    };
\end{tikzpicture}
\restoregeometry
\nopagecolor
